\documentclass[11pt]{article}
\usepackage[margin=1in]{geometry}
\usepackage{amsmath,amssymb,amsthm,bm}
\newtheorem{definition}{Definition}
\newtheorem{proposition}{Proposition}
\newcommand{\R}{\mathbb{R}}
\newcommand{\sA}{\mathcal{A}}
\newcommand{\sX}{\mathcal{X}}
\newcommand{\sL}{\mathcal{L}}
\newcommand{\ind}[1]{\mathbf{1}\!\left[#1\right]}
\newcommand{\CRP}{\mathrm{CRP}}
\newcommand{\GEM}{\mathrm{GEM}}
\newcommand{\DP}{\mathrm{DP}}
\newcommand{\Pois}{\mathrm{tPoisson}}
\title{A Hierarchical Bayesian Skill-Library Partial-Order Model}
\author{HPOP --- formal model specification}
\date{}

\begin{document}
\maketitle

\noindent
We formalize a hierarchical Bayesian model of agent traces:
\[
\textbf{Actions}\ \longrightarrow\ \textbf{Local partial orders (skills)}\ \longrightarrow\ \textbf{Global partial order},
\]
in which skills are reusable, are drawn from a shared \emph{library}, and \emph{creating a new skill is
penalized}. The model is the hierarchical, nonparametric extension of BPOP: a Dirichlet-process mixture
whose atoms are local partial orders, ordered at the top level by a global partial order. Both ordering
likelihoods use the tractable robust frontier-softmax, avoiding \#P-hard linear-extension counting.

\section{Notation}
Let $\sA=\{1,\dots,M\}$ be the canonical action (CPA) alphabet. A corpus is a set of traces; trace $n$
is an observed action sequence $y^{(n)}=(y^{(n)}_1,\dots,y^{(n)}_{T_n})$ with $y^{(n)}_t\in\sA$, obtained
after physical normalization and CPA abstraction (PDAF Steps~1a--1b). Latent structure per trace: a
segmentation of $y^{(n)}$ into contiguous \emph{skill instances}, a skill-\emph{type} label per instance,
a local partial order per skill type, and a global partial order over the instances.

\section{Partial orders and their latent representation}
\begin{definition}[Strict partial order]
A strict partial order (poset) on a finite set $\sX$ is a pair $h=(\sX,\succ_h)$ with $\succ_h$
irreflexive and transitive. $a\parallel_h b$ (\emph{incomparable}) iff neither $a\succ_h b$ nor $b\succ_h a$.
A \emph{linear extension} is a total order on $\sX$ consistent with $\succ_h$; $\sL(h)$ is the set of all
linear extensions.
\end{definition}

\begin{definition}[Latent realizer / dimension representation]
\label{def:U}
For $\sX$ with $|\sX|=m$ and a latent matrix $U\in\R^{m\times K}$, define the poset $h(U)$ by
\begin{equation}
a\succ_{h(U)} b \iff U_{a,k} > U_{b,k}\quad\forall k\in\{1,\dots,K\}.
\label{eq:dominance}
\end{equation}
Thus $h(U)$ is the intersection of $K$ total orders (\emph{realizers}); $\dim(h(U))\le K$
(Dushnik--Miller). Incomparability arises when two rows fail to dominate across all $K$ coordinates,
encoding concurrency.
\end{definition}

\noindent
\textbf{Prior on a poset.} Following Nicholls et al., place
\begin{equation}
K\sim\Pois(\lambda),\qquad U_{\cdot,k}\stackrel{iid}{\sim}\mathcal{N}(\bm 0,\Sigma_\rho),\ k=1,\dots,K,
\label{eq:Uprior}
\end{equation}
where $\Sigma_\rho$ has unit variance and correlation $\rho$; the dimension $K$ is itself inferred. Write
$p_0(h)=\int \ind{h(U)=h}\,dU\,dK$ for the induced prior over posets.

\section{The frontier-softmax likelihood}
Given a poset $h$ on $\sX$, an agent emits its elements as a stochastic linear extension by repeatedly
choosing from the \emph{frontier}.
\begin{definition}[Frontier]
For a prefix that has emitted the set $S\subseteq\sX$, the frontier is
$F(S;h)=\{a\in\sX\setminus S:\ \nexists\,b\in\sX\setminus S,\ b\succ_h a\}$,
the minimal not-yet-emitted elements. Elements of $F(S;h)$ are mutually unordered given $S$ (concurrent).
\end{definition}
\begin{definition}[Robust frontier-softmax]
\label{def:fs}
Let $w=(w_1,\dots,w_m)$ be an emitted order, $S_t=\{w_1,\dots,w_{t-1}\}$, $F_t=F(S_t;h)$,
$R_t=\sX\setminus S_t$, and the successor utility $Q(a;h)=\log\!\big(1+|\{b:a\succ_h b\}|\big)$. With
inverse temperature $\beta\ge0$ and noise $\epsilon\in[0,1)$,
\begin{equation}
p(w\mid h,\beta,\epsilon)=\prod_{t=1}^{m}\left[(1-\epsilon)\,
\frac{\ind{w_t\in F_t}\,e^{\beta Q(w_t;h)}}{\sum_{a\in F_t}e^{\beta Q(a;h)}}
\;+\;\epsilon\,\frac{1}{|R_t|}\right].
\label{eq:fs}
\end{equation}
\end{definition}
If $w_t\notin F_t$ (an order violation) the rational term vanishes and only the $\epsilon$ noise survives,
so violating traces retain positive probability. Evaluating \eqref{eq:fs} costs $O(m\,|\sX|)$ via Kahn's
algorithm, avoiding the \#P-complete enumeration of $\sL(h)$.

\section{Skills, the library, and the base measure}
A \emph{skill} is a local partial order together with its action support:
$\theta=(\sA_\theta,\,U_\theta)$ with $\sA_\theta\subseteq\sA$ and local poset $P_\theta=h(U_\theta)$ on
$\sA_\theta$. The base measure $G_0$ over skills draws a support and a poset via \eqref{eq:Uprior}.

The shared library is nonparametric:
\begin{equation}
G=\sum_{r\ge1}\pi_r\,\delta_{\theta_r},\qquad
\bm\pi\sim\GEM(\alpha),\quad \theta_r\stackrel{iid}{\sim}G_0,\qquad\text{i.e. } G\sim\DP(\alpha,G_0).
\label{eq:DP}
\end{equation}
Each skill instance's type is an i.i.d.\ draw from $G$; integrating out $\bm\pi$ yields the Chinese
restaurant process (CRP) predictive used below. (Corpus-level sharing with per-trace reuse is the
hierarchical DP; we state the single shared-library CRP, which suffices for the penalty.)

\section{Generative model}
Let $z_{n,i}\in\{1,2,\dots\}$ index the skill type of instance $i$ in trace $n$, and let
$(n,i)\mapsto g(n,i)$ enumerate all instances in exchangeable order with running type counts $n_r$.
\begin{enumerate}
\item \textbf{Library (lazy).} Maintain atoms $\theta_r=(\sA_r,U_r)$, drawing a fresh
$\theta_{r^*}\sim G_0$ the first time a new type is opened.
\item \textbf{Skill-type assignment (CRP).} For the $N{+}1$-st instance overall,
\begin{equation}
P(z=r\mid z_{1:N})=\frac{n_r}{N+\alpha},\qquad
P(z=\text{new})=\frac{\alpha}{N+\alpha}.
\label{eq:crp}
\end{equation}
\item \textbf{Global poset.} For trace $n$ with $T_n$ instances draw $V_n\in\R^{T_n\times K^g_n}$ via
\eqref{eq:Uprior} (params $\lambda_g$), giving $Q_n=h(V_n)$.
\item \textbf{Global linear extension.} Emit the instance order $\sigma_n\in\sL(Q_n)$ from the global
frontier-softmax $p(\sigma_n\mid Q_n,\beta_g,\epsilon_g)$ \eqref{eq:fs}.
\item \textbf{Local linear extensions.} For each instance $i$ of type $r=z_{n,i}$, emit its action block
$y^{(n)}_{[i]}$ (an ordering of the nodes of $P_r$) from the local frontier-softmax
$p(y^{(n)}_{[i]}\mid P_r,\beta_\ell,\epsilon_\ell)$ \eqref{eq:fs}.
\end{enumerate}
The observed trace $y^{(n)}$ is the concatenation of the blocks $\{y^{(n)}_{[i]}\}$ in the order
$\sigma_n$: a \emph{nested} linear extension (a linear extension of the global poset whose elements are
themselves linear extensions of local posets).

\section{Joint density and the new-skill penalty}
With $\Theta=\{\theta_r\}$ the realized library and $\Phi=(\alpha,\lambda,\lambda_g,\beta_\ell,\beta_g,\epsilon)$,
\begin{align}
p\big(\{y^{(n)},\sigma_n,V_n,z_{n,\cdot}\}_n,\Theta\mid\Phi\big)
=\;&\underbrace{p(z_{\cdot}\mid\alpha)}_{\text{types}}\;
\prod_{r}\underbrace{G_0(\theta_r)}_{\text{skill params}}
\prod_{n}\Big[\,\underbrace{p_0(V_n)}_{\text{global prior}}\;
\underbrace{p(\sigma_n\mid h(V_n),\beta_g,\epsilon_g)}_{\text{global LE}}\nonumber\\
&\times\prod_{i=1}^{T_n}\underbrace{p\big(y^{(n)}_{[i]}\mid P_{z_{n,i}},\beta_\ell,\epsilon_\ell\big)}_{\text{local LE}}\Big].
\label{eq:joint}
\end{align}
The first two factors are the cost of the \emph{library}; the bracket is the cost of \emph{ordering}
skills and actions. We now isolate the penalty.

\begin{proposition}[New-skill penalty via the EPPF]
\label{prop:penalty}
Let $z_{1:N}$ have $K$ distinct types with sizes $n_1,\dots,n_K$. Under the CRP \eqref{eq:crp}, the type
prior is the Ewens exchangeable partition probability function
\begin{equation}
p(z_{1:N}\mid\alpha)=\alpha^{K}\,\frac{\Gamma(\alpha)}{\Gamma(N+\alpha)}\prod_{r=1}^{K}\Gamma(n_r),
\qquad
\log p(z_{1:N}\mid\alpha)= K\log\alpha+\sum_{r}\log\Gamma(n_r)+\log\tfrac{\Gamma(\alpha)}{\Gamma(N+\alpha)} .
\label{eq:eppf}
\end{equation}
Hence the number of distinct skills $K$ enters the log-joint linearly with coefficient $\log\alpha$:
\emph{each new skill incurs a fixed prior penalty of $-\log\alpha$} (for $\alpha<1$, a positive cost),
while reusing a type of size $n_r$ contributes the cheaper $\log\Gamma(n_r{+}1)-\log\Gamma(n_r)=\log n_r$.
Opening a new type additionally pays the marginal (Bayesian-Occam) cost
\begin{equation}
-\log\!\int p\big(y_{[i]}\mid h(U)\big)\,dG_0(U)
\label{eq:occam}
\end{equation}
in place of reusing an existing $P_r$, so a new skill is favored only when it explains its block
substantially better than any library skill.
\end{proposition}

\noindent
\textbf{Pitman--Yor generalization.} Replacing the CRP by a Pitman--Yor process $\mathrm{PY}(\alpha,d)$,
$d\in[0,1)$, gives
\begin{equation}
P(z=\text{new})=\frac{\alpha+dK}{N+\alpha},\qquad P(z=r)=\frac{n_r-d}{N+\alpha},
\label{eq:py}
\end{equation}
a per-new-skill penalty $-\log(\alpha+dK)$ that grows with $K$, yielding power-law library sizes when a
heavier tail of rare skills is expected. The single knob $\alpha$ (with optional $d$) controls the
reuse--novelty trade-off: smaller $\alpha\Rightarrow$ stronger reuse and fewer skills.

\section{Posterior and inference}
The posterior over $\big(\text{segmentation},z,\{U_r\},\{V_n\},\Phi\big)$ given $\{y^{(n)}\}$ is
proportional to \eqref{eq:joint}. We target it with MCMC:
\begin{itemize}\setlength\itemsep{1pt}
\item \emph{Type reassignment} (collapsed Gibbs, DP-mixture form): for instance $(n,i)$,
$P(z_{n,i}=r\mid\cdots)\propto \big(n_r^{-(n,i)}\ \text{or}\ \alpha\big)\cdot
p\big(y^{(n)}_{[i]}\mid P_r,\beta_\ell,\epsilon_\ell\big)$, with the global-LE factor enforcing
consistency with $\sigma_n$; the $\alpha$ branch is Proposition~\ref{prop:penalty}.
\item \emph{Skill split/merge} (library moves), \emph{semi-Markov} resegmentation of $y^{(n)}$.
\item \emph{Dimension} $K_r,K^g_n$ by reversible-jump MCMC (birth/death of realizer columns).
\item \emph{Continuous} updates of $U_r,V_n,\beta_\ell,\beta_g,\epsilon$ and a hyperprior on $\alpha$ by MH.
\end{itemize}
The recovered structure compiles into a frontier executor: incomparable elements run concurrently and
reused skills are not re-planned, giving the inference-time speed-up.

\section{Informative priors from PDAF annotation}
PDAF supplies weak supervision that sharpens the priors:
\begin{itemize}\setlength\itemsep{1pt}
\item \texttt{local\_orders} (typed CPA edges with confidence) define a per-pair prior
$\pi_{ab}\in\{\succ,\prec,\parallel\}$ for each skill, replacing the uniform $1/3$ prior over the three
relations; entered as a prior factor / pseudo-observations on $U_r$. \texttt{INCOMPARABLE} contributes
direct $\parallel$ evidence, the relation hardest to recover from linear extensions alone.
\item \texttt{global\_orders} (skill-level edges) prior on $Q_n$ via $V_n$.
\item \texttt{skill\_name} and the library seed the atoms $\{\theta_r\}$ and bias the CRP labels $z$;
the annotator's \texttt{skill\_is\_new} flag is the empirical counterpart of the CRP ``new table'' event
\eqref{eq:crp} and is used as a proposal bias for opening new skills.
\end{itemize}

\end{document}
